\documentclass[a4paper,12pt,oneside]{memoir}

% Castellano
\usepackage[spanish,es-tabla]{babel}
\selectlanguage{spanish}
\usepackage[utf8]{inputenc}
\usepackage[T1]{fontenc}
\usepackage{lmodern} % Scalable font
\usepackage{microtype}
\usepackage{placeins}

\RequirePackage{booktabs}
\RequirePackage[table]{xcolor}
\RequirePackage{xtab}
\RequirePackage{multirow}

% Links
\PassOptionsToPackage{hyphens}{url}\usepackage[colorlinks]{hyperref}
\hypersetup{
	allcolors = {red}
}

% Ecuaciones
\usepackage{amsmath}

% Rutas de fichero / paquete
\newcommand{\ruta}[1]{{\sffamily #1}}

% Párrafos
\nonzeroparskip

% Huérfanas y viudas
\widowpenalty100000
\clubpenalty100000

% Imágenes

% Comando para insertar una imagen en un lugar concreto.
% Los parámetros son:
% 1 --> Ruta absoluta/relativa de la figura
% 2 --> Texto a pie de figura
% 3 --> Tamaño en tanto por uno relativo al ancho de página
\usepackage{graphicx}
\newcommand{\imagen}[3]{
	\begin{figure}[!h]
		\centering
		\includegraphics[width=#3\textwidth]{#1}
		\caption{#2}\label{fig:#1}
	\end{figure}
	\FloatBarrier
}

% Comando para insertar una imagen sin posición.
% Los parámetros son:
% 1 --> Ruta absoluta/relativa de la figura
% 2 --> Texto a pie de figura
% 3 --> Tamaño en tanto por uno relativo al ancho de página
\newcommand{\imagenflotante}[3]{
	\begin{figure}
		\centering
		\includegraphics[width=#3\textwidth]{#1}
		\caption{#2}\label{fig:#1}
	\end{figure}
}

% El comando \figura nos permite insertar figuras comodamente, y utilizando
% siempre el mismo formato. Los parametros son:
% 1 --> Porcentaje del ancho de página que ocupará la figura (de 0 a 1)
% 2 --> Fichero de la imagen
% 3 --> Texto a pie de imagen
% 4 --> Etiqueta (label) para referencias
% 5 --> Opciones que queramos pasarle al \includegraphics
% 6 --> Opciones de posicionamiento a pasarle a \begin{figure}
\newcommand{\figuraConPosicion}[6]{%
  \setlength{\anchoFloat}{#1\textwidth}%
  \addtolength{\anchoFloat}{-4\fboxsep}%
  \setlength{\anchoFigura}{\anchoFloat}%
  \begin{figure}[#6]
    \begin{center}%
      \Ovalbox{%
        \begin{minipage}{\anchoFloat}%
          \begin{center}%
            \includegraphics[width=\anchoFigura,#5]{#2}%
            \caption{#3}%
            \label{#4}%
          \end{center}%
        \end{minipage}
      }%
    \end{center}%
  \end{figure}%
}

%
% Comando para incluir imágenes en formato apaisado (sin marco).
\newcommand{\figuraApaisadaSinMarco}[5]{%
  \begin{figure}%
    \begin{center}%
    \includegraphics[angle=90,height=#1\textheight,#5]{#2}%
    \caption{#3}%
    \label{#4}%
    \end{center}%
  \end{figure}%
}
% Para las tablas
\newcommand{\otoprule}{\midrule [\heavyrulewidth]}
%
% Nuevo comando para tablas pequeñas (menos de una página).
\newcommand{\tablaSmall}[5]{%
 \begin{table}
  \begin{center}
   \rowcolors {2}{gray!35}{}
   \begin{tabular}{#2}
    \toprule
    #4
    \otoprule
    #5
    \bottomrule
   \end{tabular}
   \caption{#1}
   \label{tabla:#3}
  \end{center}
 \end{table}
}

%
% Nuevo comando para tablas pequeñas (menos de una página).
\newcommand{\tablaSmallSinColores}[5]{%
 \begin{table}[H]
  \begin{center}
   \begin{tabular}{#2}
    \toprule
    #4
    \otoprule
    #5
    \bottomrule
   \end{tabular}
   \caption{#1}
   \label{tabla:#3}
  \end{center}
 \end{table}
}

\newcommand{\tablaApaisadaSmall}[5]{%
\begin{landscape}
  \begin{table}
   \begin{center}
    \rowcolors {2}{gray!35}{}
    \begin{tabular}{#2}
     \toprule
     #4
     \otoprule
     #5
     \bottomrule
    \end{tabular}
    \caption{#1}
    \label{tabla:#3}
   \end{center}
  \end{table}
\end{landscape}
}

%
% Nuevo comando para tablas grandes con cabecera y filas alternas coloreadas en gris.
\newcommand{\tabla}[6]{%
  \begin{center}
    \tablefirsthead{
      \toprule
      #5
      \otoprule
    }
    \tablehead{
      \multicolumn{#3}{l}{\small\sl continúa desde la página anterior}\\
      \toprule
      #5
      \otoprule
    }
    \tabletail{
      \hline
      \multicolumn{#3}{r}{\small\sl continúa en la página siguiente}\\
    }
    \tablelasttail{
      \hline
    }
    \bottomcaption{#1}
    \rowcolors {2}{gray!35}{}
    \begin{xtabular}{#2}
      #6
      \bottomrule
    \end{xtabular}
    \label{tabla:#4}
  \end{center}
}

%
% Nuevo comando para tablas grandes con cabecera.
\newcommand{\tablaSinColores}[6]{%
  \begin{center}
    \tablefirsthead{
      \toprule
      #5
      \otoprule
    }
    \tablehead{
      \multicolumn{#3}{l}{\small\sl continúa desde la página anterior}\\
      \toprule
      #5
      \otoprule
    }
    \tabletail{
      \hline
      \multicolumn{#3}{r}{\small\sl continúa en la página siguiente}\\
    }
    \tablelasttail{
      \hline
    }
    \bottomcaption{#1}
    \begin{xtabular}{#2}
      #6
      \bottomrule
    \end{xtabular}
    \label{tabla:#4}
  \end{center}
}

%
% Nuevo comando para tablas grandes sin cabecera.
\newcommand{\tablaSinCabecera}[5]{%
  \begin{center}
    \tablefirsthead{
      \toprule
    }
    \tablehead{
      \multicolumn{#3}{l}{\small\sl continúa desde la página anterior}\\
      \hline
    }
    \tabletail{
      \hline
      \multicolumn{#3}{r}{\small\sl continúa en la página siguiente}\\
    }
    \tablelasttail{
      \hline
    }
    \bottomcaption{#1}
  \begin{xtabular}{#2}
    #5
   \bottomrule
  \end{xtabular}
  \label{tabla:#4}
  \end{center}
}



\definecolor{cgoLight}{HTML}{EEEEEE}
\definecolor{cgoExtralight}{HTML}{FFFFFF}

%
% Nuevo comando para tablas grandes sin cabecera.
\newcommand{\tablaSinCabeceraConBandas}[5]{%
  \begin{center}
    \tablefirsthead{
      \toprule
    }
    \tablehead{
      \multicolumn{#3}{l}{\small\sl continúa desde la página anterior}\\
      \hline
    }
    \tabletail{
      \hline
      \multicolumn{#3}{r}{\small\sl continúa en la página siguiente}\\
    }
    \tablelasttail{
      \hline
    }
    \bottomcaption{#1}
    \rowcolors[]{1}{cgoExtralight}{cgoLight}

  \begin{xtabular}{#2}
    #5
   \bottomrule
  \end{xtabular}
  \label{tabla:#4}
  \end{center}
}



\graphicspath{ {./img/} }

% Capítulos
\chapterstyle{bianchi}
\newcommand{\capitulo}[2]{
	\setcounter{chapter}{#1}
	\setcounter{section}{0}
	\setcounter{figure}{0}
	\setcounter{table}{0}
	\chapter*{\thechapter.\enskip #2}
	\addcontentsline{toc}{chapter}{\thechapter.\enskip #2}
	\markboth{#2}{#2}
}

% Apéndices
\renewcommand{\appendixname}{Apéndice}
\renewcommand*\cftappendixname{\appendixname}

\newcommand{\apendice}[1]{
	%\renewcommand{\thechapter}{A}
	\chapter{#1}
}

\renewcommand*\cftappendixname{\appendixname\ }

% Formato de portada
\makeatletter
\usepackage{xcolor}
\newcommand{\tutor}[1]{\def\@tutor{#1}}
\newcommand{\course}[1]{\def\@course{#1}}
\definecolor{cpardoBox}{HTML}{E6E6FF}
\def\maketitle{
  \null
  \thispagestyle{empty}
  % Cabecera ----------------
\noindent\includegraphics[width=\textwidth]{cabecera}\vspace{1cm}%
  \vfill
  % Título proyecto y escudo informática ----------------
  \colorbox{cpardoBox}{%
    \begin{minipage}{.8\textwidth}
      \vspace{.5cm}\Large
      \begin{center}
      \textbf{TFG del Grado en Ingeniería Informática}\vspace{.6cm}\\
      \textbf{\LARGE\@title{}}
      \end{center}
      \vspace{.2cm}
    \end{minipage}

  }%
  \hfill\begin{minipage}{.20\textwidth}
    \includegraphics[width=\textwidth]{escudoInfor}
  \end{minipage}
  \vfill
  % Datos de alumno, curso y tutores ------------------
  \begin{center}%
  {%
    \noindent\LARGE
    Presentado por \@author{}\\ 
    en Universidad de Burgos --- \@date{}\\
    Tutor: \@tutor{}\\
  }%
  \end{center}%
  \null
  \cleardoublepage
  }
\makeatother

\newcommand{\nombre}{Nombre del alumno} %%% cambio de comando

% Datos de portada
\title{título del TFG}
\author{\nombre}
\tutor{nombre tutor}
\date{\today}

\begin{document}

\maketitle


\newpage\null\thispagestyle{empty}\newpage


%%%%%%%%%%%%%%%%%%%%%%%%%%%%%%%%%%%%%%%%%%%%%%%%%%%%%%%%%%%%%%%%%%%%%%%%%%%%%%%%%%%%%%%%
\thispagestyle{empty}


\noindent\includegraphics[width=\textwidth]{cabecera}\vspace{1cm}

\noindent D. nombre tutor, profesor del departamento de nombre departamento, área de nombre área.

\noindent Expone:

\noindent Que el alumno D. \nombre, con DNI dni, ha realizado el Trabajo final de Grado en Ingeniería Informática titulado título de TFG. 

\noindent Y que dicho trabajo ha sido realizado por el alumno bajo la dirección del que suscribe, en virtud de lo cual se autoriza su presentación y defensa.

\begin{center} %\large
En Burgos, {\large \today}
\end{center}

\vfill\vfill\vfill

% Author and supervisor
\begin{minipage}{0.45\textwidth}
\begin{flushleft} %\large
Vº. Bº. del Tutor:\\[2cm]
D. nombre tutor
\end{flushleft}
\end{minipage}
\hfill
\begin{minipage}{0.45\textwidth}
\begin{flushleft} %\large
Vº. Bº. del co-tutor:\\[2cm]
D. nombre co-tutor
\end{flushleft}
\end{minipage}
\hfill

\vfill

% para casos con solo un tutor comentar lo anterior
% y descomentar lo siguiente
%Vº. Bº. del Tutor:\\[2cm]
%D. nombre tutor


\newpage\null\thispagestyle{empty}\newpage




\frontmatter

% Abstract en castellano
\renewcommand*\abstractname{Resumen}
\begin{abstract}
En este primer apartado se hace una \textbf{breve} presentación del tema que se aborda en el proyecto.
\end{abstract}

\renewcommand*\abstractname{Descriptores}
\begin{abstract}
Palabras separadas por comas que identifiquen el contenido del proyecto Ej: servidor web, buscador de vuelos, android \ldots
\end{abstract}

\clearpage

% Abstract en inglés
\renewcommand*\abstractname{Abstract}
\begin{abstract}
A \textbf{brief} presentation of the topic addressed in the project.
\end{abstract}

\renewcommand*\abstractname{Keywords}
\begin{abstract}
keywords separated by commas.
\end{abstract}

\clearpage

% Indices
\tableofcontents

\clearpage

\listoffigures

\clearpage

\listoftables
\clearpage

\mainmatter
\capitulo{1}{Introducción}

\section{Marco del trabajo}

El turismo es uian actividad que conlleva que las personas se desplacen
de su lugar habitual de forma temporal para visitar otros lugares, ya sea por
ocio, negocios, cultura u otras razones. Esta actividad puede ser entendida
como una forma de consumo ya que los viajeros gastan dinero en bienes y
servicios durante su estancia en el destino turístico.\cite{turismo}

En España, el turismo es una de las principales fuentes de ingresos y
empleo. En 2023, la actividad turística supuso un 13,4\% del PIB del país.
Esto no es casualidad, sus playas, gastronomía, clima mediterráneo y la
riqueza artística y arquitectónica han logrado posicionarlo entre los tres
países más visitados en todo el mundo llegando a recibir alrededor de 85
millones de turistas en ese mismo año.\cite{turismoEspana2025}.

Debido a la importancia del turismo en la economía española, la Sociedad
Mercantil Estatal para la Gestión de la Innovación y las Tecnologías Turísticas (SEGITTUR) lleva varios años trabajando para contribuir al desarrollo,
modernización y mantenimento de la industria turística española a través
de la innovación tecnológica. Trata de mejorar la competividad, calidad
y sostenibilidad en los ámbitos medioambientales, económicos y sociales
relacionados con el turismo.\cite{segitturEjes}

Con el objetivo de mejorar la competividad de los destinos turísticos
y la calidad de vida de los viajeros, la Secretaría de Estado de Turismo
(SETUR) promueve el programa Destino Turístico Inteligente (DTI) el cual
es gestionado por SEGITTUR. Un DTI es un destino turístico innovador con
una accesibilidad global y que cuenta con una infraestructura tecnológica avanzada. Esta infraestructura garantiza el desarrollo sostenible del destino
y trata de mejorar la experiencia de los viajeros y su calidad de vida.\cite{segittur:dti}

\section{Contenido del trabajo}

El objetivo principal del trabajo es la creacón y publicación de un cuadro
de mando con el objetivo de que la provincia de Burgos alcance el distintivo
de Destino Turístico Inteligente. Para ello, se ha desarrollado un sistema
de extracción automatizado de reseñas de puntos de interés en la provincia
de Burgos. Con esa información se crea el cuadro de mando que permite
analizar y monitorizar distintas métricas relevantes como la distribución de
los puntos de interés por categorias. Además, se ha desarrollado una red
neuronal entrenada con un conjunto de datos formado por reseñas extraídas
de la misma forma con el objetivo de clasificar por categorías las reseñas
proporcionadas externamente.

\section{Estructura de la memoria}

La memoria del proyecto se estructura en siete capítulos:

\begin{itemize}
    \item \textbf{Capítulo 1: Introducción.} Presenta el marco del trabajo, el contenido y la estructura de la memoria.
    \item \textbf{Capítulo 2: Objetivos del proyecto.} Define los objetivos generales y específicos del proyecto.
    \item \textbf{Capítulo 3: Conceptos teóricos.} : Explica los conceptos teóricos necesarios para entender el proyecto.
    \item \textbf{Capítulo 4: Técnicas y herramientas.} Describe las técnicas y herramientas utilizadas en el desarrollo del proyecto.
    \item \textbf{Capítulo 5: Aspectos relevantes del desarrollo del proyecto.} Detalla el desarrollo del proyecto, incluyendo la extracción de reseñas, la creación del cuadro de mando y el entrenamiento de la red neuronal.
    \item \textbf{Capítulo 6: Trabajos relacionados.} Presenta una revisión de trabajos relacionados con el proyecto, incluyendo estudios previos y proyectos similares.
    \item \textbf{Capítulo 7: Conclusiones y líneas de trabajo futuras.} Resume las conclusiones del proyecto y propone posibles trabajos futuros.
\end{itemize}

Además, se incluten varios anexos que contienen información adicional relevante para el proyecto:

\begin{itemize}
    \item \textbf{Anexo A: Plan de proyecto software.} Contiene el plan de proyecto, incluyendo la planificación temporal y el desarrollo del mismo y la viabilidad económica y legal.
    \item \textbf{Anexo B: Especificación de requisitos.} Incluye los requisitos funcionales y no funcionales del proyecto, así como los casos de uso. 
    \item \textbf{Anexo C: Especificación de diseño.} Explica el diseño del sistema, incluyendo la arquitectura o las interfaces entre otras cosas.
    \item \textbf{Anexo D: Documentación técnica de programación.} Describe la parte técnica del proyecto como las explicaciones sobre el repositorio del código fuente.
    \item \textbf{Anexo E: Documentación de usuario.} Incluye de forma detallada todo lo que el usuario necesita saber y poseer para poder ejecutar y/o utilizar el proyecto principalmente en forma de manual de usuario.
    \item \textbf{Anexo F: Anexo de sostenibilización curricular.} Presenta una reflexión sobre los aspectos de sostenibilidad abordados en el proyecto.
\end{itemize}

\section{Materiales entregados}

Los materiales entregados al finalizar el proyecto son:

\begin{itemize}
    \item Varios scripts de Python utilizados para la extracción de reseñas.
    \item Un conjunto de datos con las reseñas extraídas utilizadas para entrenar la red neuronal.
    \item Un modelo de red neuronal entrenado para clasificar reseñas por categorías.   
    \item Un cuadro de mando interactivo creado con Power BI publicado en Power Pages.
    \item Una memoria del proyecto junto a sus anexos en formato PDF creado con \LaTeX.
    \item Un repositorio de GitHub con el código fuente del proyecto y la documentación.\cite{repositorio}
\end{itemize}

\capitulo{2}{Objetivos del proyecto}

Este proyecto tiene como finalidad desarrollar una solución tecnológica que permita analizar la reputación online de distintos destinos turísticos.

\section {Crear un cuadro de mando}
Crear un cuadro de mando destinado al análisis de la reputación online, aplicable a distintos destinos turísticos.
Este cuadro de mando contendrá información visual y estructurada sobre la percepción de los usuarios sobre dichos destinos turísticos.

\section {Sistema de análisis de sentimientos}
El análisis de sentimientos permite interpretar de forma automática la opinión de los usuarios sobre un destino turístico.
Se desarrollará un sistema de análisis de sentimientos que permita obtener información sobre la percepción de los usuarios.
Se utilizarán reseñas de usuarios para extraer dicha información sobre los destinos turísticos.

\section {Clasificar e interpretar los datos}
Hacer una clasficación e interpretación de los datos obtenidos: recursos, reseñas, etc. El objetivo de este proceso es obtener información relevante.
Para ello, se clasificarán las reseñas de los usuarios en positivas, negativas o neutras.
También se aplicara el mismo procedimiento para las palabras más frecuentes en las reseñas.
Esta tarea se realizará con PowerBI, una herramienta de análisis de datos.

\section {Automatizar la obtención de datos}
 
Para que el sistema de análisis de reputación online sea eficiente y aplicable a múltiples destinos turísticos, es fundamental automatizar el proceso de recopilación de datos. 
Esto permitirá que el sistema escale sin intervención manual constante, garantizando una actualización en tiempo real y una mayor precisión en la información analizada.

\section {Mostrar los resultados}
Mostrar dichos resultados a través de una página web creada con PowerPages con el fin de que los datos sean accesibles y comprensibles fácilmente por los usuarios.
Dicha plataforma permite la creación de páginas web de forma sencilla y rápida sin necesidad de escribir código.

\capitulo{3}{Propuestas de extracción de datos}

\section{Propuestas de extracción de datos}
\begin{itemize}
    \item \textbf{TripAdvisor API:} La API de TripAdvisor contiene millones de reviews de usuarios sobre destinos turísticos de todo el mundo. Debido a sus límites de reviews (5 por lugar, es decir, 5 reviews por solicitud) no es factible.
    Ofrecen 5000 peticiones gratuitas al mes. 
    \begin{table}[h!]
        \centering
        \begin{tabular}{|c|c|}
            \hline
            \textbf{Solicitudes} & \textbf{Costo por solicitud} \\
            \hline
            0 - 5,000 & €0.00 \\
            5,001 - 20,000 & €0.00924 \\
            20,001 - 100,000 & €0.00859 \\
            100,001 - 500,000 & €0.00804 \\
            500,000+ & €0.00758 \\
            \hline
        \end{tabular}
        \caption{Costos por solicitud en la API de Tripadvisor}
    \end{table}
    \item \textbf{API de Google Places (textSearch) y Google Maps Reviews Scraper de Apify:} Utilizando textSearch se puede obtener un identificador para cada POI de forma gratuita e ilimitada. Para ello, es necesario tener un punto de referencia (coordenadas) y un radio de búsqueda.  
    Los puntos de referencia se pueden obtener mediante una consulta en Overpass Query Language desde su interfaz web para extraer las coordenadas de OSM. Al contrario que en el caso anterior, ahora no se emplean los puntos límites de los municipios, solo un punto céntrico de cada uno de ellos.
    El problema en este caso es que no se obtienen todos los POI de un municipio, sino solo los que se encuentran dentro del radio de búsqueda.
    \item \textbf{API de Google Places (nearbySearch) y ficheros binarios de OpenStreetMaps(OSM):} OSM proporciona ficheros binarios (.osm.pbf) que se pueden convertir a GeoJSON para obtener información relevante sobre los municipios como las coordenadas de sus límites. 
    Con la API de Google Places se pueden extraer reviews utilizando nearbySearch que permite obtener las reviews de los POI más relevantes que se encuentren cerca de un punto dado por sus coordenadas. 
    Sin embargo, esta API solo permite obtener 5 reviews por lugar y 60 lugares por cada búsqueda alrededor de un punto. Esto implica que si se quiere obtener información de un municipio con más de 60 POI, se deben realizar múltiples búsquedas y muchas peticiones.
    \item \textbf{Consultas de Overpass QL, API de Google Places (nearbySearch) y Google Maps Reviews Scraper de Apify:} A través de la interfaz web de Overpass QL
    Al igual que en el caso anterior, se pueden obtener los límites de los municipios mediante consultas en Overpass QL.
    Esta herramienta permite extraer reviews de Google Maps de forma ilimitada y gratuita.
    Con estos identificadores se puede automatizar la extracción de reviews a través de Apify que ofrece 10000 reviews gratuitas por cada cuenta creada.
    También ofrecen un descuento del 50\% para cuentas educativas. Las dos opciones más interesantes son la cuenta Profesional (18,06€/mes tras el descuento del 50\%) que permite extraer 78000 reviews al mes y la cuenta Scale (92,16€/mes tras el descuento del 50\%) que permite extraer 398000 reviews al mes. También ofrecen un descuento de 10\% si se paga anualmente.
\end{itemize}
\capitulo{4}{Conceptos teóricos}

En aquellos proyectos que necesiten para su comprensión y desarrollo de unos conceptos teóricos de una determinada materia o de un determinado dominio de conocimiento, debe existir un apartado que sintetice dichos conceptos.

Algunos conceptos teóricos de \LaTeX{} \footnote{Créditos a los proyectos de Álvaro López Cantero: Configurador de Presupuestos y Roberto Izquierdo Amo: PLQuiz}.

\section{Secciones}

Las secciones se incluyen con el comando section.

\subsection{Subsecciones}

Además de secciones tenemos subsecciones.

\subsubsection{Subsubsecciones}

Y subsecciones. 


\section{Referencias}

Las referencias se incluyen en el texto usando cite~\cite{wiki:latex}. Para citar webs, artículos o libros~\cite{koza92}, si se desean citar más de uno en el mismo lugar~\cite{bortolot2005, koza92}.


\section{Imágenes}

Se pueden incluir imágenes con los comandos standard de \LaTeX, pero esta plantilla dispone de comandos propios como por ejemplo el siguiente:

\imagen{escudoInfor}{Autómata para una expresión vacía}{.5}



\section{Listas de items}

Existen tres posibilidades:

\begin{itemize}
	\item primer item.
	\item segundo item.
\end{itemize}

\begin{enumerate}
	\item primer item.
	\item segundo item.
\end{enumerate}

\begin{description}
	\item[Primer item] más información sobre el primer item.
	\item[Segundo item] más información sobre el segundo item.
\end{description}
	
\begin{itemize}
\item 
\end{itemize}

\section{Tablas}

Igualmente se pueden usar los comandos específicos de \LaTeX o bien usar alguno de los comandos de la plantilla.

\tablaSmall{Herramientas y tecnologías utilizadas en cada parte del proyecto}{l c c c c}{herramientasportipodeuso}
{ \multicolumn{1}{l}{Herramientas} & App AngularJS & API REST & BD & Memoria \\}{ 
HTML5 & X & & &\\
CSS3 & X & & &\\
BOOTSTRAP & X & & &\\
JavaScript & X & & &\\
AngularJS & X & & &\\
Bower & X & & &\\
PHP & & X & &\\
Karma + Jasmine & X & & &\\
Slim framework & & X & &\\
Idiorm & & X & &\\
Composer & & X & &\\
JSON & X & X & &\\
PhpStorm & X & X & &\\
MySQL & & & X &\\
PhpMyAdmin & & & X &\\
Git + BitBucket & X & X & X & X\\
Mik\TeX{} & & & & X\\
\TeX{}Maker & & & & X\\
Astah & & & & X\\
Balsamiq Mockups & X & & &\\
VersionOne & X & X & X & X\\
} 

\capitulo{5}{Técnicas y herramientas}

Esta parte de la memoria tiene como objetivo presentar las técnicas metodológicas y las herramientas de desarrollo que se han utilizado para llevar a cabo el proyecto. Si se han estudiado diferentes alternativas de metodologías, herramientas, bibliotecas se puede hacer un resumen de los aspectos más destacados de cada alternativa, incluyendo comparativas entre las distintas opciones y una justificación de las elecciones realizadas. 
No se pretende que este apartado se convierta en un capítulo de un libro dedicado a cada una de las alternativas, sino comentar los aspectos más destacados de cada opción, con un repaso somero a los fundamentos esenciales y referencias bibliográficas para que el lector pueda ampliar su conocimiento sobre el tema.



\capitulo{6}{Aspectos relevantes del desarrollo del proyecto}

Este apartado pretende recoger los aspectos más interesantes del desarrollo del proyecto, comentados por los autores del mismo.
Debe incluir desde la exposición del ciclo de vida utilizado, hasta los detalles de mayor relevancia de las fases de análisis, diseño e implementación.
Se busca que no sea una mera operación de copiar y pegar diagramas y extractos del código fuente, sino que realmente se justifiquen los caminos de solución que se han tomado, especialmente aquellos que no sean triviales.
Puede ser el lugar más adecuado para documentar los aspectos más interesantes del diseño y de la implementación, con un mayor hincapié en aspectos tales como el tipo de arquitectura elegido, los índices de las tablas de la base de datos, normalización y desnormalización, distribución en ficheros3, reglas de negocio dentro de las bases de datos (EDVHV GH GDWRV DFWLYDV), aspectos de desarrollo relacionados con el WWW...
Este apartado, debe convertirse en el resumen de la experiencia práctica del proyecto, y por sí mismo justifica que la memoria se convierta en un documento útil, fuente de referencia para los autores, los tutores y futuros alumnos.

\capitulo{7}{Trabajos relacionados}

Este apartado sería parecido a un estado del arte de una tesis o tesina. En un trabajo final grado no parece obligada su presencia, aunque se puede dejar a juicio del tutor el incluir un pequeño resumen comentado de los trabajos y proyectos ya realizados en el campo del proyecto en curso. 

\capitulo{8}{Conclusiones y Líneas de trabajo futuras}

Todo proyecto debe incluir las conclusiones que se derivan de su desarrollo. Éstas pueden ser de diferente índole, dependiendo de la tipología del proyecto, pero normalmente van a estar presentes un conjunto de conclusiones relacionadas con los resultados del proyecto y un conjunto de conclusiones técnicas. 
Además, resulta muy útil realizar un informe crítico indicando cómo se puede mejorar el proyecto, o cómo se puede continuar trabajando en la línea del proyecto realizado. 



\bibliographystyle{plain}
\bibliography{bibliografia}

\end{document}
